% ============================================================
\chapter{Word Embeddings}
\label{ch:embeddings}
% ============================================================

In 2013, Tomas Mikolov and his team at Google published a paper that would transform natural language processing: \emph{Efficient Estimation of Word Representations in Vector Space}. The idea was to train a shallow neural network to predict words from their context (or vice versa) and use the learned weights as vector representations of words. The result, Word2Vec, revealed a stunning geometric structure: semantic analogies become \emph{vector operations}. The vector ``king'' $-$ ``man'' $+$ ``woman'' points to ``queen.'' Geometry captures meaning.

But Word2Vec did not emerge from nothing. It belongs to a long tradition founded on the \emph{distributional hypothesis} of J.R.\ Firth (1957): ``you shall know a word by the company it keeps.'' This chapter traces the path from co-occurrence matrices to dense embeddings, through latent semantic analysis and GloVe.

\begin{intuition}[title=\textbf{The distributional hypothesis}]
``You shall know a word by the company it keeps'' (J.R.\ Firth, 1957). Word
embeddings are dense vector representations that capture semantic similarity:
words appearing in similar contexts have nearby vectors.
\end{intuition}

% ------------------------------------------------------------
\section{From Sparse to Dense Representations}
\label{sec:emb:motivation}
% ------------------------------------------------------------

Recall that a one-hot representation encodes each word $w_i$ as a vector
$\mathbf{e}_i \in \R^V$ with $(\mathbf{e}_i)_j = \delta_{ij}$. This
representation has two major drawbacks:
\begin{enumerate}
    \item \textbf{Dimensionality}: $V$ can reach $10^5$ to $10^6$.
    \item \textbf{Orthogonality}: $\inner{\mathbf{e}_i}{\mathbf{e}_j} = 0$
          for $i \neq j$, even for synonyms.
\end{enumerate}

The goal of embeddings is to learn a function
$\phi : \mathcal{V} \to \R^d$ with $d \ll V$ (typically $d \in [100, 300]$)
such that geometric proximity reflects semantic similarity.

% ------------------------------------------------------------
\section{Word2Vec}
\label{sec:emb:word2vec}
% ------------------------------------------------------------

\subsection{CBOW Architecture}

\begin{definition}[Continuous Bag of Words (CBOW)]
The CBOW model predicts the center word $w_t$ from its context
$C_t = \{w_{t-c}, \ldots, w_{t-1}, w_{t+1}, \ldots, w_{t+c}\}$ where $c$ is
the window size. The objective is to maximize:
\[
\mathcal{L}_{\text{CBOW}} = \sum_{t=1}^{T} \log P(w_t \mid C_t)
\]
with:
\[
P(w_t \mid C_t) = \frac{\exp(\inner{\mathbf{v}_{w_t}}{\bar{\mathbf{u}}_{C_t}})}{\sum_{w \in \mathcal{V}} \exp(\inner{\mathbf{v}_w}{\bar{\mathbf{u}}_{C_t}})}
\]
where $\bar{\mathbf{u}}_{C_t} = \frac{1}{2c} \sum_{w \in C_t} \mathbf{u}_w$
is the average of context vectors.
\end{definition}

\subsection{Skip-gram Architecture}

\begin{definition}[Skip-gram]
The Skip-gram model predicts context words from the center word. The objective
is to maximize:
\[
\mathcal{L}_{\text{SG}} = \sum_{t=1}^{T} \sum_{\substack{j=-c \\ j \neq 0}}^{c} \log P(w_{t+j} \mid w_t)
\]
with:
\[
P(w_o \mid w_i) = \frac{\exp(\inner{\mathbf{v}_{w_o}}{\mathbf{u}_{w_i}})}{\sum_{w \in \mathcal{V}} \exp(\inner{\mathbf{v}_w}{\mathbf{u}_{w_i}})}
\]
\end{definition}

\begin{warning}[title=\textbf{Softmax cost}]
The softmax denominator requires a sum over the entire vocabulary $V$, which
is prohibitive. Two common approximations are used: negative sampling and
hierarchical softmax.
\end{warning}

\subsection{Negative Sampling (NEG)}

\begin{definition}[Negative Sampling]
The negative sampling objective replaces the softmax with:
\[
\mathcal{L}_{\text{NEG}} = \log \sigma(\inner{\mathbf{v}_{w_o}}{\mathbf{u}_{w_i}}) + \sum_{k=1}^{K} \E_{w_k \sim P_n} \left[\log \sigma(-\inner{\mathbf{v}_{w_k}}{\mathbf{u}_{w_i}})\right]
\]
where $K$ is the number of negative examples and $P_n(w) \propto \text{count}(w)^{3/4}$.
\end{definition}

\begin{theorem}[NEG-PMI Equivalence]
\label{thm:neg-pmi}
Levy and Goldberg (2014) showed that Skip-gram with negative sampling
implicitly factorizes a shifted PMI matrix:
\[
\mathbf{u}_w^\top \mathbf{v}_c = \text{PMI}(w, c) - \log K
\]
where $\text{PMI}(w, c) = \log \frac{P(w, c)}{P(w) P(c)}$ is the pointwise
mutual information.
\end{theorem}

% ------------------------------------------------------------
\section{GloVe}
\label{sec:emb:glove}
% ------------------------------------------------------------

\begin{definition}[GloVe]
The GloVe model (Pennington et al., 2014) learns embeddings by minimizing the
weighted squared error on the co-occurrence matrix:
\[
\mathcal{L}_{\text{GloVe}} = \sum_{i,j=1}^{V} f(X_{ij}) \left(\inner{\mathbf{u}_i}{\mathbf{v}_j} + b_i + c_j - \log X_{ij}\right)^2
\]
where $X_{ij}$ is the co-occurrence count, $b_i, c_j$ are biases, and $f$ is a
weighting function:
\[
f(x) = \begin{cases}
(x / x_{\max})^\alpha & \text{if } x < x_{\max} \\
1 & \text{otherwise}
\end{cases}
\]
with typically $x_{\max} = 100$ and $\alpha = 0.75$.
\end{definition}

\begin{remark}
GloVe combines the advantages of global methods (matrix factorization) and
local methods (context window). The GloVe objective is equivalent to a
factorization of the log-co-occurrence matrix.
\end{remark}

% ------------------------------------------------------------
\section{FastText}
\label{sec:emb:fasttext}
% ------------------------------------------------------------

\begin{definition}[FastText]
FastText (Bojanowski et al., 2017) enriches Skip-gram by representing each word
as the sum of its character $n$-grams. For a word $w$ with $n$-gram set
$\mathcal{G}_w$:
\[
\mathbf{u}_w = \sum_{g \in \mathcal{G}_w} \mathbf{z}_g
\]
where $\mathbf{z}_g \in \R^d$ is the vector of $n$-gram $g$.
\end{definition}

\begin{bestpractice}[title=\textbf{FastText advantages}]
\begin{itemize}
    \item Handles out-of-vocabulary (OOV) words via $n$-gram decomposition.
    \item Captures morphology (prefixes, suffixes).
    \item Particularly effective for morphologically rich languages (Turkish,
          Finnish, Arabic).
\end{itemize}
\end{bestpractice}

% ------------------------------------------------------------
\section{Geometric Properties}
\label{sec:emb:geometry}
% ------------------------------------------------------------

\subsection{Vector Analogies}

One of the most remarkable properties of embeddings is their ability to capture
semantic relations through arithmetic operations:

\begin{example}
\[
\mathbf{v}_{\text{king}} - \mathbf{v}_{\text{man}} + \mathbf{v}_{\text{woman}} \approx \mathbf{v}_{\text{queen}}
\]
\end{example}

Formally, to solve the analogy $a : b :: c : ?$, we seek:
\[
d^* = \argmax_{d \in \mathcal{V} \setminus \{a,b,c\}} \cos(\mathbf{v}_d, \mathbf{v}_b - \mathbf{v}_a + \mathbf{v}_c)
\]

\begin{theorem}[Linear Structure of Analogies]
Under the assumption that embeddings factorize a PMI matrix
(Theorem~\ref{thm:neg-pmi}), regular semantic relations manifest as linear
directions in embedding space:
\[
\mathbf{v}_b - \mathbf{v}_a \approx \mathbf{v}_{b'} - \mathbf{v}_{a'}
\]
for all pairs $(a, b)$ and $(a', b')$ sharing the same relation.
\end{theorem}

\subsection{t-SNE Visualization}

\begin{codeblock}[title=\textbf{Embedding visualization}]
\begin{minted}{python}
import numpy as np
from sklearn.manifold import TSNE
import matplotlib.pyplot as plt
from gensim.models import KeyedVectors

# Load pre-trained embeddings
wv = KeyedVectors.load_word2vec_format('GoogleNews-vectors.bin', binary=True)
words = ['king', 'queen', 'man', 'woman', 'prince', 'princess',
         'paris', 'france', 'berlin', 'germany', 'rome', 'italy']
vectors = np.array([wv[w] for w in words])

tsne = TSNE(n_components=2, random_state=42, perplexity=5)
coords = tsne.fit_transform(vectors)

plt.figure(figsize=(10, 8))
for i, word in enumerate(words):
    plt.scatter(coords[i, 0], coords[i, 1])
    plt.annotate(word, (coords[i, 0]+0.5, coords[i, 1]+0.5))
plt.title("t-SNE Visualization of Word2Vec Embeddings")
plt.show()
\end{minted}
\end{codeblock}

% ------------------------------------------------------------
\section{Embedding Evaluation}
\label{sec:emb:evaluation}
% ------------------------------------------------------------

\subsection{Intrinsic Evaluation}

\begin{itemize}
    \item \textbf{Analogies}: accuracy on analogy benchmarks (Google analogy
          dataset, BATS).
    \item \textbf{Similarity}: Spearman correlation between cosine similarity
          and human judgments (SimLex-999, WordSim-353).
\end{itemize}

\subsection{Extrinsic Evaluation}

Performance of embeddings as features in downstream tasks: sentiment
classification, NER, word sense disambiguation.

\begin{keyformulas}[title=\textbf{Embedding evaluation metrics}]
\begin{align}
\text{Analogy accuracy} &= \frac{\text{correct analogies}}{\text{total analogies}} \\
\rho_{\text{Spearman}} &= 1 - \frac{6 \sum_i d_i^2}{n(n^2 - 1)} & \text{(rank correlation)}
\end{align}
\end{keyformulas}

% ------------------------------------------------------------
\section{Bias in Embeddings}
\label{sec:emb:bias}
% ------------------------------------------------------------

Embeddings learned from large text corpora inherit societal biases present in
the data.

\begin{example}
Bolukbasi et al.\ (2016) showed that:
\[
\mathbf{v}_{\text{computer programmer}} - \mathbf{v}_{\text{man}} + \mathbf{v}_{\text{woman}} \approx \mathbf{v}_{\text{homemaker}}
\]
This result reflects gender stereotypes present in the training corpora.
\end{example}

\begin{definition}[Debiasing by Projection]
The debiasing method of Bolukbasi et al.\ consists of:
\begin{enumerate}
    \item Identifying the bias direction $\mathbf{g}$ (via PCA on gender pairs).
    \item Projecting neutral words orthogonally to $\mathbf{g}$:
          $\mathbf{v}_w' = \mathbf{v}_w - \inner{\mathbf{v}_w}{\mathbf{g}} \mathbf{g}$.
    \item Equalizing definitional pairs.
\end{enumerate}
\end{definition}

% ------------------------------------------------------------
\section{Contextual vs.\ Static Embeddings}
\label{sec:emb:contextual}
% ------------------------------------------------------------

Static embeddings (Word2Vec, GloVe, FastText) assign a single vector to each
word, regardless of context. This is problematic for polysemous words:

\begin{example}
The word ``bank'' can refer to a financial institution or the side of a river.
A static embedding represents both senses with a single vector, which is a
compromise between the two.
\end{example}

Contextual embeddings (ELMo, BERT, GPT) solve this problem by producing
different representations depending on context. They will be studied in
Chapter~\ref{ch:pretrained}.

% ------------------------------------------------------------
\section{Complete Implementation}
\label{sec:emb:implementation}
% ------------------------------------------------------------

\begin{codeblock}[title=\textbf{Skip-gram training with PyTorch}]
\begin{minted}{python}
import torch
import torch.nn as nn

class SkipGram(nn.Module):
    def __init__(self, vocab_size, embed_dim):
        super().__init__()
        self.target_embed = nn.Embedding(vocab_size, embed_dim)
        self.context_embed = nn.Embedding(vocab_size, embed_dim)

    def forward(self, target, context, negatives):
        # target: (B,), context: (B,), negatives: (B, K)
        t = self.target_embed(target)          # (B, d)
        c = self.context_embed(context)        # (B, d)
        n = self.context_embed(negatives)      # (B, K, d)

        # Positive score
        pos_score = torch.sum(t * c, dim=1)    # (B,)
        pos_loss = -torch.log(torch.sigmoid(pos_score) + 1e-10)

        # Negative scores
        neg_score = torch.bmm(n, t.unsqueeze(2)).squeeze(2)  # (B, K)
        neg_loss = -torch.log(torch.sigmoid(-neg_score) + 1e-10).sum(1)

        return (pos_loss + neg_loss).mean()

model = SkipGram(vocab_size=50000, embed_dim=300)
optimizer = torch.optim.Adam(model.parameters(), lr=0.001)
\end{minted}
\end{codeblock}

% ------------------------------------------------------------
\section{Exercises}
\label{sec:emb:exercises}
% ------------------------------------------------------------

\begin{exercise}[$\star$]
Show that the cosine similarity between two one-hot vectors is zero for
different words.
\end{exercise}

\begin{exercise}[$\star\star$]
Derive the gradient of the Skip-gram objective with negative sampling with
respect to vectors $\mathbf{u}_{w_i}$ and $\mathbf{v}_{w_o}$.
\end{exercise}

\begin{exercise}[$\star\star$]
Show that the GloVe objective is equivalent to a weighted matrix factorization
of the matrix $\log X_{ij}$.
\end{exercise}

\begin{exercise}[$\star\star$]
Train Word2Vec embeddings (using \texttt{gensim}) on an English corpus
(Wikipedia). Evaluate the analogies ``king -- man + woman = ?'' and
``Paris -- France + Germany = ?''.
\end{exercise}

\begin{exercise}[$\star\star\star$]
Implement the debiasing method of Bolukbasi et al.\ (2016). Measure gender
bias before and after correction on a set of profession-related words.
\end{exercise}
